




IRS 63 is a Class I protostar embedded in the L1709 region of the Ophiuchus molecular cloud at a distance of 144 pc \citep{sadavoy_dust_2019, Dom2020}. It is one of the brightest Class I protostars at (sub)millimetre wavelengths, and is surrounded by a relatively large circumstellar disk with a radius larger than 50 AU \citep{Dom2020}. 



Since IRS 63 is relatively bright and nearby, it is an ideal target for polarization observations. Additionally, the large disk allows the polarization morphology to be studied from the inner to outer disk. Lastly, since IRS 63 is a Class I protostar, it provides the opportunity to study dust grain growth during the early stages of planet formation \citep{Dom2020}. 



\add{previous observations and what they found}